
\documentclass{article}
\usepackage{colortbl}
\usepackage{makecell}
\usepackage{multirow}
\usepackage{supertabular}

\begin{document}

\newcounter{utterance}

\centering \large Interaction Transcript for game `cladder', experiment `full\_v1.5\_default', episode 4908 with score{-}Qwen3.5{-}9B{-}sp.
\vspace{24pt}

{ \footnotesize  \setcounter{utterance}{1}
\setlength{\tabcolsep}{0pt}
\begin{supertabular}{c@{$\;$}|p{.15\linewidth}@{}p{.15\linewidth}p{.15\linewidth}p{.15\linewidth}p{.15\linewidth}p{.15\linewidth}}
    \# & \multicolumn{2}{c}{Player} && \multicolumn{2}{c}{Game Master} \\
    \hline

    \theutterance \stepcounter{utterance}  
    & & & \multicolumn{4}{p{0.6\linewidth}}{
        \cellcolor[rgb]{0.9,0.9,0.9}{
            \makecell[{{p{\linewidth}}}]{
                \texttt{\tiny{[P1$\langle$GM]}}
                \texttt{You are an expert in causal inference. The following question is not a typical commonsense query, but rather a meticulously designed question created by a professor specializing in causal inference, intended to assess the students' mastery of the course content.} \\
\texttt{Imagine a self{-}contained, hypothetical world with only the following conditions, and without any unmentioned factors or causal relationships: My decision has a direct effect on penguin mood and curly hair. Penguin mood has a direct effect on curly hair. For those who choose to take the stairs, the probability of curly hair is 70\%. For those who choose to take the elevator, the probability of curly hair is 17\%. For those who choose to take the elevator, would it be less likely to see curly hair if I had taken the stairs?} \\
\texttt{Start your answer with 'Yes' or 'No', followed by additional reasoning or evidence to support your explanation.} \\
            }
        }
    }
     \\ \\

    \theutterance \stepcounter{utterance}  
    & \multicolumn{4}{p{0.6\linewidth}}{
        \cellcolor[rgb]{0.561,0.737,0.561}{
            \makecell[{{p{\linewidth}}}]{
                \texttt{\tiny{[P1$\rangle$GM]}}
                \texttt{No} \\
\\ 
\texttt{**Reasoning:**} \\
\\ 
\texttt{To determine the answer, we must analyze the causal structure and the specific counterfactual query using the provided conditions.} \\
\\ 
\texttt{1.  **Causal Structure**:} \\
\texttt{    *   **Decision** \$\textbackslash{}rightarrow\$ **Penguin Mood**} \\
\texttt{    *   **Decision** \$\textbackslash{}rightarrow\$ **Curly Hair**} \\
\texttt{    *   **Penguin Mood** \$\textbackslash{}rightarrow\$ **Curly Hair**} \\
\texttt{    *   This implies that "Decision" is a common cause of both "Penguin Mood" and "Curly Hair," and "Penguin Mood" is a mediator between "Decision" and "Curly Hair."} \\
\\ 
\texttt{2.  **The Counterfactual Query**:} \\
\texttt{    The question asks: "For those who choose to take the elevator, would it be less likely to see curly hair if I had taken the stairs?"} \\
\texttt{    *   **Factual World (\$W\_\{actual\}\$)**: Decision = Elevator.} \\
\texttt{    *   **Counterfactual World (\$W\_\{counterfactual\}\$)**: Decision = Stairs.} \\
\texttt{    *   **Target Variable**: Probability of Curly Hair (\$P(\textbackslash{}text\{Curly Hair\})\$).} \\
\texttt{    *   **Comparison**: Is \$P(\textbackslash{}text\{Curly Hair\} | \textbackslash{}text\{Stairs\}) < P(\textbackslash{}text\{Curly Hair\} | \textbackslash{}text\{Elevator\})\$?} \\
\\ 
\texttt{3.  **Evaluating the Probabilities**:} \\
\texttt{    *} \\
            }
        }
    }
    & & \\ \\

    \theutterance \stepcounter{utterance}  
    & & & \multicolumn{2}{p{0.3\linewidth}}{
        \cellcolor[rgb]{0.9,0.9,0.9}{
            \makecell[{{p{\linewidth}}}]{
                \texttt{\tiny{[GM$|$GM]}}
                \texttt{No} \\
            }
        }
    }
    & & \\ \\

    \theutterance \stepcounter{utterance}  
    & & & \multicolumn{2}{p{0.3\linewidth}}{
        \cellcolor[rgb]{0.9,0.9,0.9}{
            \makecell[{{p{\linewidth}}}]{
                \texttt{\tiny{[GM$|$GM]}}
                \texttt{no} \\
            }
        }
    }
    & & \\ \\

    \theutterance \stepcounter{utterance}  
    & & & \multicolumn{2}{p{0.3\linewidth}}{
        \cellcolor[rgb]{0.9,0.9,0.9}{
            \makecell[{{p{\linewidth}}}]{
                \texttt{\tiny{[GM$|$GM]}}
                \texttt{game\_result = WIN} \\
            }
        }
    }
    & & \\ \\

\end{supertabular}
}

\end{document}
